Show for cuda how 
Show an illsutration of how cuda architectuer works the streaming multiprocessors and how we have parallelixewd in in two diffretn ways. 
Then possibly show the speeed gain betwen chunk vs normal along wiht a little snippet that shows how we call the cude kernel and how it is implemented 
using teh block and thread ids. Also that we scale using the streaming multirocessors.
Just take the figure from the cuda book


Maybe compare the two methods where chunks gets differnt chunks sizes??

this si our two implementations



__device__ typedef struct MatrixChunk {
    int row_start;
    int row_end;
    int col_start;
    int col_end;
} MatrixChunk;

__device__ MatrixChunk chunk_from_index(int i, int rows, int cols, int chunk_size) {
    int row = i * chunk_size / rows * chunk_size;
    int col = i * chunk_size % rows;
    int row_end = rows - row >= chunk_size ? row + chunk_size : rows;
    int col_end = cols - col >= chunk_size ? col + chunk_size : cols;

    return (MatrixChunk) {
        .row_start = row,
        .row_end = row_end,
        .col_start = col,
        .col_end = col_end
    };
}



__global__ void matmul_cuda_transposed(Matrix a, Matrix bT, Matrix c) {
    int total_threads = blockDim.x * gridDim.x;

    for (int workIndex = threadIdx.x + blockDim.x * blockIdx.x;
         workIndex < a.rows * a.cols;
         workIndex += total_threads)
    {
        int i = workIndex / a.cols;
        int j = workIndex % a.cols;

        float sum = 0.0f;

        for (int k = 0; k < a.cols; k++) {
            float a_val = a.ptr[i * a.cols + k];
            float b_val = bT.ptr[j * bT.cols + k];
            sum += a_val * b_val;
        }

        c.ptr[i * c.cols + j] = sum;
    }
}

__global__ void matmul_cuda_chunks(Matrix a, Matrix b, Matrix c, int chunk_size) {
    int total_threads = blockDim.x * gridDim.x;
    int work_index = threadIdx.x + blockDim.x * blockIdx.x;
    int chunk_count = ((c.rows - 1) / chunk_size + 1) * ((c.cols - 1) / chunk_size + 1);

    for (int chunk_index = work_index; chunk_index < chunk_count; chunk_index += total_threads) {
        MatrixChunk chunk = chunk_from_index(chunk_index, c.rows, c.cols, chunk_size);
        for (int i = chunk.row_start; i < chunk.row_end; i++) {
            for (int k = 0; k < a.cols; k++) {
                float entry = a.ptr[i * a.cols+ k]; //*matrix_index(a, i, k);
                for (int j = chunk.col_start; j < chunk.col_end; j++) {
                    c.ptr[i * c.cols + j] += entry * b.ptr[k * b.cols + j]; // *matrix_index(b, k, j);
                }
            }
        }
    }
}


I was wondering whoch diagrams i should take, maybe one that shoes differne between cpu and gpu and rhe oen fir Streamin multirpocseeors to understand those.
Then hiw much if these two implementations ant the additonla MatrixChunk code should i inlcude in slides

Maybe also for diffrent matrix dimensions show the abegre runtime for normal cude and for chunk cuda with chunk size 4. 